\section{Experiments}

\subsection{Experimental Setup}

\subsubsection{Datasets}

We evaluate DVCL on four heterogeneous graph benchmarks: ACM, DBLP, AMiner, and IMDB. Each dataset contains multiple node and relation types, and the task is semi-supervised classification of a target node type. Following the HSeCo protocol, dataset-specific meta-paths are used to construct semantic neighborhoods. Table~\ref{tab:dataset_settings} summarizes the target node type and meta-paths. The final main comparison will contain at least three datasets with complete five-seed results; the fourth dataset will be reported when the same protocol has been completed.

\begin{table}[htbp]
\centering
\caption{Dataset settings.}
\label{tab:dataset_settings}
\begin{tabular}{lll}
\toprule
Dataset & Target node type & Meta-paths \\
\midrule
ACM & Paper & P-A-P, P-F-P \\
DBLP & Author & A-P-A, A-P-C-P-A, A-P-T-P-A \\
AMiner & Paper & P-A-P, P-R-P \\
IMDB & Movie & M-D-M, M-A-M \\
\bottomrule
\end{tabular}
\end{table}

\subsubsection{Evaluation Protocol}

We use classification accuracy as the only evaluation metric:
\begin{equation}
\operatorname{Accuracy}
=\frac{1}{|\mathcal{V}_{\mathrm{eval}}|}
\sum_{i\in\mathcal{V}_{\mathrm{eval}}}
\mathbb{I}(\hat{y}_i=y_i).
\end{equation}
Unless otherwise stated, each table reports the mean and standard deviation over seeds 1, 2, 3, 4, and 5. All methods use the same train/validation/test split, attacked graph artifact, target-node IDs, and evaluation nodes for each seed. Hyperparameters are selected only on the validation split under an equal tuning budget. A dash (``--'') denotes a result that has not yet been produced under the unified protocol. When a paired HSeCo reference is available, the small colored line below an accuracy value reports its change relative to the HSeCo row in the same column; green and ``+'' indicate an increase, whereas red and ``-'' indicate a decrease. Preliminary tables use the explicitly labeled pre-audit HSeCo row, while final tables will use HSeCo$^{\dagger}$.

Numbers copied from the original papers are not mixed with rerun results in the main tables. HSeCo$^{\dagger}$ denotes our audited reproduction. The difference between published HSeCo results and our reproduction will be documented separately in the appendix, together with the data split, configuration, and checkpoint behavior.

\subsubsection{Baselines}

Table~\ref{tab:baselines} defines the baseline set. It covers standard heterogeneous architectures, heterogeneous contrastive learning, and defenses designed for structural perturbations. Homogeneous defenses such as GCN-Jaccard and GNNGuard are excluded from the primary comparison unless the same pre-specified graph conversion is applied to every such method.

\begin{table}[htbp]
\centering
\small
\caption{Methods included in the unified comparison.}
\label{tab:baselines}
\begin{tabular}{lp{0.23\linewidth}p{0.56\linewidth}}
\toprule
Method & Category & Role in the comparison \\
\midrule
HAN~\cite{wang2019han} & Heterogeneous backbone & Meta-path attention baseline without an explicit robustness mechanism. \\
HGT~\cite{hu2020hgt} & Heterogeneous backbone & Type-aware Transformer baseline using typed projections and attention. \\
MAGNN~\cite{fu2020magnn} & Heterogeneous backbone & Meta-path instance aggregation baseline with fine-grained semantic modeling. \\
HeCo~\cite{wang2021heco} & Contrastive HGNN & Heterogeneous contrastive-learning baseline without an explicit attack defense. \\
HSeCo$^{\dagger}$~\cite{hseco} & Robust HGNN & Direct semantic-topology baseline, rerun by us under the unified protocol. \\
RoHe~\cite{rohe} & Robust HGNN & Heterogeneous defense baseline designed for adversarial structural perturbations. \\
HeteGuard~\cite{heteguard} & Robust HGNN & Heterogeneous defense baseline that suppresses unreliable structural information. \\
DVCL & Proposed method & Purified semantic topology, feature-induced graph, and cross-view contrastive alignment. \\
\bottomrule
\end{tabular}
\end{table}

\subsubsection{Attack Settings}

For global attacks, we evaluate PRBCD and HetePRBCD with perturbation rates of $5\%$, $10\%$, $15\%$, $20\%$, and $25\%$. PRBCD performs global structural perturbation, whereas HetePRBCD preserves heterogeneous relation constraints during attack generation. Clean accuracy is reported alongside attacked accuracy.

For target attacks, we use an FGSM-style heterogeneous structural attack. On ACM, 500 target paper nodes are selected from the test split for each seed, and the same target IDs are reused by all methods. The final target-attack table reports budgets $\Delta\in\{1,2,3\}$ and evaluates accuracy only on the selected target nodes.

The current attack suite modifies graph topology while leaving node attributes unchanged. Because DVCL constructs one view from node features, we additionally reserve experiments for feature perturbation and a joint adaptive attack that optimizes against both DVCL views. Until those experiments are complete, the empirical claim is restricted to robustness against structural perturbations.

\subsubsection{Implementation and Model Selection}

The maximum number of training epochs is 200, early stopping is based on validation performance, and Adam is used for optimization. Each checkpoint contains every trainable component required for inference. For HSeCo$^{\dagger}$, this includes both the HAN semantic module and the node classifier; after checkpoint restoration, the semantic graph is regenerated from the restored module.

DVCL uses raw target-node features, directed KNN, and $k=20$ for the feature-induced view. The topology view follows the HSeCo semantic topology construction, and the two view-specific representations are concatenated for classification. We set $\lambda_{\mathrm{HAN}}=1.0$. Based on the preliminary sweep in Table~\ref{tab:param}, we fix $\lambda_{\mathrm{DVCL}}=1.0$ before the unified rerun. This setting balances clean and attacked accuracy, and it will not be retuned on individual attack test sets or datasets.

\FloatBarrier
\subsection{Main Results}

\subsubsection{Cross-Dataset Summary}

Table~\ref{tab:cross_dataset} is the final summary table to be completed after the unified rerun. ``Attack Avg.'' is the arithmetic mean over the ten PRBCD and HetePRBCD conditions. The currently populated ACM entries are preliminary records; they are retained to guide the paper draft but are not treated as the final external-baseline comparison.

\begin{table}[htbp]
\centering
\small
\setlength{\tabcolsep}{3.5pt}
\caption{Cross-dataset accuracy under the unified protocol. Final results will be reported as mean $\pm$ standard deviation over five seeds.}
\label{tab:cross_dataset}
\resizebox{\textwidth}{!}{%
\begin{tabular}{lcccccccc}
\toprule
& \multicolumn{2}{c}{ACM} & \multicolumn{2}{c}{DBLP} & \multicolumn{2}{c}{AMiner} & \multicolumn{2}{c}{IMDB} \\
\cmidrule(lr){2-3}\cmidrule(lr){4-5}\cmidrule(lr){6-7}\cmidrule(lr){8-9}
Method & Clean & Attack Avg. & Clean & Attack Avg. & Clean & Attack Avg. & Clean & Attack Avg. \\
\midrule
HAN & -- & -- & -- & -- & -- & -- & -- & -- \\
HGT & -- & -- & -- & -- & -- & -- & -- & -- \\
MAGNN & -- & -- & -- & -- & -- & -- & -- & -- \\
HeCo & -- & -- & -- & -- & -- & -- & -- & -- \\
HSeCo$^{\dagger}$ & -- & -- & -- & -- & -- & -- & -- & -- \\
RoHe & -- & -- & -- & -- & -- & -- & -- & -- \\
HeteGuard & -- & -- & -- & -- & -- & -- & -- & -- \\
DVCL ($\lambda=1.0$, prelim.) & $0.8810\pm0.0049$ & 0.8723 & -- & -- & -- & -- & -- & -- \\
DVCL ($\lambda=1.0$, final) & -- & -- & -- & -- & -- & -- & -- & -- \\
\bottomrule
\end{tabular}}
\end{table}

\subsubsection{Global Structural Attacks on ACM}

Tables~\ref{tab:acm_prbcd} and~\ref{tab:acm_heteprbcd} use models as rows so that every baseline can be added without changing the reporting format. The pre-audit HSeCo row and the DVCL ($\lambda=0.5$) row contain available historical records; their colored changes are paired within each column. They will be replaced by HSeCo$^{\dagger}$ and the fixed $\lambda=1.0$ DVCL configuration after all methods have been run on the same frozen artifacts.

\begin{table}[htbp]
\centering
\small
\caption{Accuracy on ACM under PRBCD. Values currently present are preliminary five-seed records.}
\label{tab:acm_prbcd}
\resizebox{\textwidth}{!}{%
\begin{tabular}{lcccccc}
\toprule
Method & Clean & 5\% & 10\% & 15\% & 20\% & 25\% \\
\midrule
HAN & -- & -- & -- & -- & -- & -- \\
HGT & -- & -- & -- & -- & -- & -- \\
MAGNN & -- & -- & -- & -- & -- & -- \\
HeCo & -- & -- & -- & -- & -- & -- \\
HSeCo (pre-audit) & $0.8679\pm0.0129$ & $0.8491\pm0.0086$ & $0.8519\pm0.0080$ & $0.8484\pm0.0091$ & $0.8530\pm0.0092$ & $0.8504\pm0.0092$ \\
HSeCo$^{\dagger}$ (final) & -- & -- & -- & -- & -- & -- \\
RoHe & -- & -- & -- & -- & -- & -- \\
HeteGuard & -- & -- & -- & -- & -- & -- \\
DVCL ($\lambda=0.5$, prelim.) & \accgain{$0.8770\pm0.0071$}{+0.91 pp} & \accgain{$0.8732\pm0.0092$}{+2.41 pp} & \accgain{$0.8718\pm0.0077$}{+1.99 pp} & \accgain{$0.8674\pm0.0060$}{+1.90 pp} & \accgain{$0.8660\pm0.0146$}{+1.30 pp} & \accgain{$0.8659\pm0.0082$}{+1.55 pp} \\
DVCL ($\lambda=1.0$, final) & -- & -- & -- & -- & -- & -- \\
\bottomrule
\end{tabular}}
\end{table}

\begin{table}[htbp]
\centering
\small
\caption{Accuracy on ACM under HetePRBCD. Values currently present are preliminary five-seed records.}
\label{tab:acm_heteprbcd}
\resizebox{\textwidth}{!}{%
\begin{tabular}{lcccccc}
\toprule
Method & Clean & 5\% & 10\% & 15\% & 20\% & 25\% \\
\midrule
HAN & -- & -- & -- & -- & -- & -- \\
HGT & -- & -- & -- & -- & -- & -- \\
MAGNN & -- & -- & -- & -- & -- & -- \\
HeCo & -- & -- & -- & -- & -- & -- \\
HSeCo (pre-audit) & $0.8679\pm0.0129$ & $0.8530\pm0.0121$ & $0.8589\pm0.0126$ & $0.8546\pm0.0070$ & $0.8548\pm0.0124$ & $0.8511\pm0.0092$ \\
HSeCo$^{\dagger}$ (final) & -- & -- & -- & -- & -- & -- \\
RoHe & -- & -- & -- & -- & -- & -- \\
HeteGuard & -- & -- & -- & -- & -- & -- \\
DVCL ($\lambda=0.5$, prelim.) & \accgain{$0.8770\pm0.0071$}{+0.91 pp} & \accgain{$0.8762\pm0.0027$}{+2.33 pp} & \accgain{$0.8737\pm0.0087$}{+1.48 pp} & \accgain{$0.8726\pm0.0046$}{+1.79 pp} & \accgain{$0.8700\pm0.0108$}{+1.53 pp} & \accgain{$0.8703\pm0.0063$}{+1.92 pp} \\
DVCL ($\lambda=1.0$, final) & -- & -- & -- & -- & -- & -- \\
\bottomrule
\end{tabular}}
\end{table}

The available ACM records indicate that the preliminary DVCL configuration is more stable than the pre-audit HSeCo reproduction across the tested structural perturbation rates. This observation remains provisional: the final claim will be based only on post-audit results produced by the frozen protocol and the fixed $\lambda_{\mathrm{DVCL}}=1.0$ configuration.

\subsubsection{Target Structural Attacks}

Table~\ref{tab:target_attack} reports target-node accuracy at three attack budgets. Existing records are available only at $\Delta=3$ for HSeCo and the preliminary DVCL configuration. Budgets 1 and 2, external baselines, and the final DVCL configuration remain to be completed on the same target-node IDs.

\begin{table}[htbp]
\centering
\small
\caption{Target-node accuracy on ACM under the FGSM-style structural attack.}
\label{tab:target_attack}
\begin{tabular}{lcccc}
\toprule
Method & Clean & $\Delta=1$ & $\Delta=2$ & $\Delta=3$ \\
\midrule
HAN & -- & -- & -- & -- \\
HGT & -- & -- & -- & -- \\
MAGNN & -- & -- & -- & -- \\
HeCo & -- & -- & -- & -- \\
HSeCo (pre-audit) & 0.8624 & -- & -- & 0.7608 \\
HSeCo$^{\dagger}$ (final) & -- & -- & -- & -- \\
RoHe & -- & -- & -- & -- \\
HeteGuard & -- & -- & -- & -- \\
DVCL ($\lambda=0.5$, prelim.) & \accgain{0.8724}{+1.00 pp} & -- & -- & \accgain{0.8608}{+10.00 pp} \\
DVCL ($\lambda=1.0$, final) & -- & -- & -- & -- \\
\bottomrule
\end{tabular}
\end{table}

\FloatBarrier
\subsection{Ablation and Sensitivity Analysis}

\subsubsection{Contrastive-Loss Weight}

Table~\ref{tab:param} summarizes the available sweep of $\lambda_{\mathrm{DVCL}}$. The attack average covers all ten global attack conditions. We select $\lambda_{\mathrm{DVCL}}=1.0$ because it gives the best clean accuracy and nearly the strongest attacked accuracy, providing a balanced fixed configuration for the unified rerun. The value 2.0 is not selected solely for its slightly higher attack average.

\begin{table}[htbp]
\centering
\caption{Accuracy sensitivity to $\lambda_{\mathrm{DVCL}}$ on ACM.}
\label{tab:param}
\begin{tabular}{ccc}
\toprule
$\lambda_{\mathrm{DVCL}}$ & Clean Accuracy & Attack Average Accuracy \\
\midrule
0.1 & -- & 0.8665 \\
0.5 & $0.8770\pm0.0071$ & 0.8710 \\
1.0 & $0.8810\pm0.0049$ & 0.8723 \\
2.0 & $0.8790\pm0.0085$ & 0.8734 \\
\bottomrule
\end{tabular}
\end{table}

\subsubsection{Component Ablation}

Table~\ref{tab:component_ablation} separates the effects of the two views and cross-view contrastive learning. These preliminary records use the raw directed KNN view with $k=20$ and the earlier $\lambda_{\mathrm{DVCL}}=0.5$ setting; the final ablation will use the same frozen split and artifacts as the main comparison.

\begin{table}[htbp]
\centering
\caption{Component ablation accuracy on ACM.}
\label{tab:component_ablation}
\begin{tabular}{lcc}
\toprule
Variant & Clean Accuracy & Attack Average Accuracy \\
\midrule
Full DVCL & $0.8770\pm0.0071$ & 0.8710 \\
DVCL w/o contrastive loss & $0.8743\pm0.0074$ & 0.8656 \\
Topology view only & $0.8523\pm0.0177$ & 0.8471 \\
Feature view only & $0.8539\pm0.0086$ & 0.8539 \\
\bottomrule
\end{tabular}
\end{table}

\subsubsection{Feature-View Construction}

Table~\ref{tab:knn_supp} reports a preliminary three-seed comparison of feature-view constructions. This analysis is supplementary and does not change the fixed raw directed KNN configuration used in the main comparison.

\begin{table}[htbp]
\centering
\caption{Feature-view construction accuracy on ACM over three seeds.}
\label{tab:knn_supp}
\begin{tabular}{lcc}
\toprule
Feature view & Clean Accuracy & Attack Average Accuracy \\
\midrule
Raw KNN, $k=10$ & 0.8750 & 0.8719 \\
Raw KNN, $k=20$ & 0.8752 & 0.8706 \\
Raw KNN, $k=40$ & 0.8665 & 0.8677 \\
Mutual raw KNN, $k=20$ & 0.8804 & 0.8739 \\
Heterogeneous-profile KNN, $k=20$ & 0.8807 & 0.8784 \\
\bottomrule
\end{tabular}
\end{table}

\FloatBarrier
\subsection{Additional Robustness Analysis}

Table~\ref{tab:adaptive_attack} reserves the tests needed to determine whether the feature-induced view remains reliable when the attacker is aware of the complete DVCL architecture. These results are required before making claims beyond structural robustness.

\begin{table}[htbp]
\centering
\small
\caption{Accuracy under attacks beyond topology-only perturbation.}
\label{tab:adaptive_attack}
\begin{tabular}{lcccc}
\toprule
Attack setting & HSeCo$^{\dagger}$ & RoHe & HeteGuard & DVCL \\
\midrule
Feature perturbation & -- & -- & -- & -- \\
Joint topology--feature attack & -- & -- & -- & -- \\
Adaptive dual-view attack & -- & -- & -- & -- \\
\bottomrule
\end{tabular}
\end{table}

\FloatBarrier
\subsection{Reproducibility and Claim Boundary}

Before the final comparison, we audit HSeCo$^{\dagger}$ by correcting early stopping, jointly saving and restoring the semantic module and classifier, and regenerating the semantic graph from the restored checkpoint. For every dataset and seed, we store a manifest containing the split indices, attack artifact path, target-node IDs, random seed, model configuration, and selected checkpoint. All baselines are then rerun from these frozen manifests.

At the current stage, the evidence supports only a restricted conclusion: under the unified ACM protocol, preliminary DVCL results are stronger than the current HSeCo reproduction under PRBCD, HetePRBCD, and the FGSM-style structural target attack. A broader claim against existing robust heterogeneous graph methods requires completed results for multiple datasets and the external baselines in Table~\ref{tab:baselines}.

\FloatBarrier
